% SPDX-License-Identifier: CC-BY-ND-4.0
\documentclass[
% handout,
aspectratio=169,
xcolor={usenames}
]{beamer}
% % To install handoutWithNotes from GitHub:
% % git clone https://github.com/gdiepen/latexbeamer-handoutWithNotes
% % cd latexbeamer-handoutWithNotes/
% % tlmgr install l3build
% % l3build install
% \usepackage{handoutWithNotes}
% \pgfpagesuselayout{4 on 1 with notes}[letterpaper, border shrink=5mm]
\graphicspath{ {img/} }
\setbeamertemplate{navigation symbols}{}
\setbeamertemplate{footline}[frame number]
\setbeamercolor{footline}{fg=white}
\setbeamertemplate{footline}{
  \vspace{0.02\paperheight}
  \hspace{4pt}\textbf{May 5--8, 2025 \quad Chicago, IL}\quad\textbf{\color{hpsf-green}\#HPSFcon}\hfill\insertframenumber\hspace{4pt}
  \vspace{0.02\paperheight}
}
\setbeamercolor{author}{fg=white}
\setbeamercolor{institute}{fg=white}
\setbeamercolor{date}{fg=white}
\renewcommand{\footnoterule}{}
\definecolor{hpsf-green}{RGB}{0,176,80}

\usepackage{tikz}
\usetikzlibrary{calc}           % ($ ... $)
\usetikzlibrary{decorations.pathmorphing}    % zigzag
\usetikzlibrary{fit}            % fit
\usetikzlibrary{positioning}    % of
\usetikzlibrary{shapes.symbols} % signal
\tikzset{
  invisible/.style={opacity=0,text opacity=0},
  visible on/.style={alt=#1{}{invisible}},
  alt/.code args={<#1>#2#3}{%
    \alt<#1>{\pgfkeysalso{#2}}{\pgfkeysalso{#3}}
  },
}

\addtobeamertemplate{frametitle}{%
  \begin{tikzpicture}[remember picture, overlay]
    \node (bg) [anchor = north, inner sep = 0pt, outer sep = 0pt]
    at (current page.north) {%
      \includegraphics[width=\paperwidth]{im-hpsf-bg.png}};
    \fill [white]
    (current page.north east) rectangle
    % 0.45 / 7.51 = 0.06
    ($ (current page.south west) + (0, 0.06\paperheight) $);
    \node [
    anchor = north west,
    inner sep = 0pt,
    outer sep = 0pt,
    % 11.08 / 13.33 = 0.8312078
    xshift = 0.8312078\paperwidth,
    %  0.49 /  7.51 = 0.06524634
    yshift = -0.06524634\paperheight
    ]
    at (current page.north west) {%
      % 1.93 / 13.33 = 0.1447862
      \includegraphics[width=0.1447862\paperwidth]{logo-hpsf-black.png}};
  \end{tikzpicture}
}{
}

\usepackage[
backend=biber,
style=authoryear,
sorting=none
]{biblatex}
\addbibresource{\jobname.bib}

\usepackage{soul}               % \ul

\usepackage{datetime2}          % \DTMnow

\title{Challenges mixing Spack-optimized %
  hardware accelerator libraries on %
  personal scientific computers}
\author{Pariksheet Nanda}
\institute{University of Pittsburgh\\%
  Chemical Engineering}
\date{May 8, 2025}

\begin{document}

\section{Title}

{
\setbeamercolor{titlelike}{fg=white}
% Hide frame number on the title slide using noframenumbering,plain options per
% https://tex.stackexchange.com/a/317983
\begin{frame}[noframenumbering, plain]
  \begin{tikzpicture}[remember picture, overlay]
    \node [
    anchor = north,
    inner sep = 0pt,
    outer sep = 0pt
    ]
    at (current page.north) {%
      \includegraphics[width=\paperwidth]{im-hpsf-bg.png}};
    \node [
    color = hpsf-green,
    anchor = north west,
    % 0.33 / 13.33 = 0.02475619
    xshift = +0.02475619\paperwidth,
    % 0.49 /  7.51 = 0.06524634
    yshift = -0.06524634\paperheight
    ] at (current page.north west) {\textbf{\#HPSFcon}};
    \node [
    anchor = north west,
    inner sep = 0pt,
    outer sep = 0pt,
    % 9.69 / 13.33 = 0.7269317
    % 0.49 /  7.51 = 0.06524634
    xshift = +0.726931\paperwidth,
    yshift = -0.06524634\paperheight
    ]
    at (current page.north west) {%
      % 3.32 / 13.33 = 0.2490623
      \includegraphics[width=0.2490623\paperwidth]{logo-hpsf.png}};
    % Sponsors:
    % Linux Foundation
    % DOE Advanced Scientific Computing Research Program
    % Wellcome Leap
    % NIH NHGRI
    % NIH NIAID
    % NSF ACCESS
    \node (lf) at (current page.west)
    [anchor = west, xshift = 1cm, yshift = 0.5cm] {
      \includegraphics[width=.2\framewidth]{logo-linux-foundation.pdf}};
    \node (doe) at (lf.south west) [anchor = north west, yshift = -0.2cm] {
      \includegraphics[width=.25\framewidth]{logo-doe.pdf}};
    \node (anl) at (doe.south west) [anchor = north west, yshift = -0.2cm] {
      \includegraphics[width=.2\framewidth]{logo-argonne.png}};

    \node (leap) at (current page.east)
    [anchor = east, xshift = -1cm, yshift = 0.2cm] {
      \includegraphics[width=.12\framewidth]{logo-leap.pdf}};
    \node (nhgri) at ($ (leap.south west) + (-1.5cm,0) $)
    [anchor = north west, yshift = -0.2cm] {
      \includegraphics[height=.055\framewidth]{logo-nhgri.pdf}};
    \node (niaid) at (nhgri.south west)
    [anchor = north west, yshift = -0.2cm] {
      \includegraphics[height=.06\framewidth]{logo-niaid.pdf}};

    \node (umich) at (current page.south)
    [anchor = south, yshift = 0.5cm] {
      \includegraphics[height=.035\framewidth]{logo-umich.pdf}};
    \node (uconn) at (umich.north west)
    [anchor = north east, xshift = -0.5cm] {
      \includegraphics[height=.05\framewidth]{logo-uconn.pdf}};
    \node (pitt) at (umich.east)
    [anchor = west, xshift = -0.75cm, yshift = -0.1cm] {
      \includegraphics[height=.25\framewidth]{logo-pitt.pdf}};

    \node at (current page.south east)
    [anchor = south east, text width = 2.2cm,
    execute at begin node = \setlength{\baselineskip}{6pt}] {%
      \color{white}\scriptsize\par Logos of funders of HPC or spack activities};
  \end{tikzpicture}
  \maketitle{}
\end{frame}
}

\section{Background}

\subsection{Background about myself}
\begin{frame}{Background about myself}
  \begin{tikzpicture}[
    remember picture,
    overlay,
    timepoint/.style = {
      circle,
      fill = black,
      inner sep = 1pt
    },
    description/.style ={
      inner sep = 0pt,
      text width = 4.5cm
    }
    ]
    \coordinate (y07) at ($ (current page.west) +(+1,0.7) $);
    \coordinate (y24) at ($ (current page.east) +(-1,0.7) $);
    \coordinate (y16) at ($ (y07) !.4! (y24) $);
    \pgfmathparse{(2022 - 2016) / (2024 - 2016)}
    \coordinate (y22) at ($ (y16) !\pgfmathresult! (y24) $);
    \path [draw]
    (y07)
    node [timepoint] {}
    node [below] {2007}
    node [above, anchor = south] {\includegraphics[scale=.08]{logo-gentoo.png}}
    node
    [description, below = .75 of y07, anchor = north west, xshift = -1cm] {
      \begin{itemize}[<+->]
      \item Low barrier from\\user $\to$ developer
      \item Use flags\\$\to$ spack variants
      \end{itemize}
    }
    decorate
    [decoration = {name = zigzag, pre length = 1cm, post length = 1cm}]
    { (y07) -- (y16) }
    (y16)
    node [timepoint] {}
    node [below] {2016}
    node [above, anchor = south, xshift = -1cm] {\includegraphics<+->[scale=.3]{logo-sc16.png}}
    node [description, below = .75 of y16, anchor = north] {
      \begin{itemize}[<+->]
      \item<.-> HPC supervisor tells me about spack
      \item Installed 33 packages on university cluster
      \item Set spack as default on research lab server
      \end{itemize}
    }
    --
    (y22)
    node [timepoint] {}
    node [below] {2022}
    node [above, anchor = south, xshift = -1.8cm] {%
      \includegraphics<+->[scale=0.5,clip,trim={8 8 10 5}]{im-gran.png}}
    node [description, below = .75 of y22, anchor = north, text width = 5cm,
    xshift = -1cm] {
      \begin{itemize}[<+->]
      \item<.-> Calibrate drug dosing for multiscale model of lung infection
      \item<.-> Each sample 16~hrs
      \item<.-> 2~week walltime
      \end{itemize}
    }
    --
    (y24)
    node [timepoint] {}
    node [below] {2024};
    \node at (y24) [above, text width = 1.8cm, visible on = {<+->}]
    {Wikipedia article: ``Spack (package $\phantom{(}$manager)''};
  \end{tikzpicture}
\end{frame}

\subsection{Motivation for this talk}
\begin{frame}{Motivation for this talk}
  \begin{columns}[T]
    \begin{column}{.6\framewidth}
      \begin{itemize}[<+->]
      \item Laptops now come with neural processing units (NPUs)
        \begin{description}[<.->][2024]
        \item[2023] Intel Meteor Lake AI Boost
        \item[2023] AMD Ryzen XDNA
        \item[2024] Qualcomm Snapdragon X Hexagon
        \item[2024] M-series Apple Silicon Apple Intelligence
        \end{description}
      \item Bridging GPU / NPU scientific libraries %
        from Linux to Windows and macOS
      \item R (and python): NPU hardware optimization
      \end{itemize}
    \end{column}
    \begin{column}{.4\framewidth}
      \begin{tikzpicture}[%
        remember picture,
        overlay,
        every node/.style = {
          inner sep = 0pt,
          outer sep = 0pt
        }]
        \begin{scope}
          \coordinate (intel-center) at (1.5cm, -1cm);
          % Viewport.
          \clip ([xshift = -1.5cm, yshift = 1cm]intel-center)
          rectangle +(3cm, -2cm);
          \node (intel) at (intel-center) {%
            \includegraphics[height=2cm,clip,trim={0 135 30 130}]%
            {npu-intel.png}};
        \end{scope}
        \begin{scope}
          \coordinate [xshift = 3cm + 1pt] %
          (qualcomm-center) at (intel-center);
          % Viewport.
          \clip ([xshift = -1.5cm, yshift = -1cm]qualcomm-center)
          rectangle +(3cm, 2cm);
          \node (qualcomm) at (qualcomm-center) {%
            \includegraphics[height=2cm,clip,trim={0 110 0 130}]%
            {npu-qualcomm.jpg}};
        \end{scope}
        \begin{scope}
          \coordinate [yshift = -2cm -1pt] %
          (apple-center) at (qualcomm-center);
          % Viewport.
          \clip ([xshift = -1.5cm, yshift = -1cm]apple-center)
          rectangle +(3cm, 2cm);
          \node (apple) at (apple-center) {%
            \includegraphics[height=2cm,clip,trim={0 55 0 55}]{npu-apple.jpg}};
        \end{scope}
        \begin{scope}
          \coordinate [yshift = -2cm -1pt] %
          (amd-center) at (intel-center);
          % Viewport.
          \clip ([xshift = -1.5cm, yshift = -1cm]amd-center)
          rectangle +(3cm, 2cm);
          \node (amd) at (amd-center) {%
            \includegraphics[height=2cm,clip,trim={0 170 0 170}]{npu-amd.jpg}};
        \end{scope}
      \end{tikzpicture}
    \end{column}
  \end{columns}
\end{frame}

\begin{frame}{Unlike other presentations\ldots}
  \ldots I'm \ul{proposing} ideas for early feedback,
  not reporting anything that I've done
\end{frame}

\subsection{Foreign function interfaces (FFI) in R and python}
\begin{frame}{Foreign function interfaces (FFI) in R and python}
  \begin{columns}[T]
    \begin{column}{.85\framewidth}
      \begin{itemize}[<+->]
      \item Python
        \begin{itemize}[<.->]
        \item Native C/C++ API extensions rarely used; too verbose
        \item Cython, cffi, SWIG, Numba
        \end{itemize}
      \item R
        \begin{itemize}
        \item<.-> Native extensions API (Fortran, C, C++17) easy to use%
          \footnotemark[1]
        \item<.-> Easy to debug with \texttt{R -d gdb} / \texttt{R -d lldb}
        \item<.-> Common to lightly wrap C++ libraries + headers as R packages %
          (e.g. Boost, Armadillo)
        \item Probably your favorite FFI for a higher language, %
          even if you don't know it yet :)
        \end{itemize}
      \item Mostly going to talk about R because of funding
      \end{itemize}
    \end{column}
    \begin{column}{.15\framewidth}
      \begin{tikzpicture}[%
        every node/.style = {
          inner sep = 0pt,
          outer sep = 0pt
        }
        ]
        \node (r)
        {\includegraphics[width=.5\textwidth]{logo-r.png}};
        \node [anchor = west, xshift = 10pt, yshift = -2pt] (py) at (r.east)
        {\includegraphics[width=.4\textwidth]{logo-python.png}};
        \node (lang) [%
        draw,
        fit = (r) (py),
        inner sep = 2pt,
        signal,
        signal pointer angle = 150,
        signal to = south] {};
        \path let \p1 = (lang.north west), \p2 = (lang.south east) in
        node (ffi) at (lang.south) [%
        anchor = north,
        yshift = -5pt,
        minimum width = \x2 - \x1,
        minimum height = \y1 - \y2,
        draw,
        signal,
        signal pointer angle = 150,
        signal from = north,
        signal to = south] {FFI};
        \path let \p1 = (lang.north west), \p2 = (lang.south east) in
        node (dots) at (ffi.south) [%
        anchor = north,
        yshift = -5pt,
        minimum width = \x2 - \x1,
        minimum height = \y1 - \y2,
        draw,
        signal,
        signal pointer angle = 150,
        signal from = north,
        signal to = south]
        {\ldots};
        \path let \p1 = (lang.north west), \p2 = (lang.south east) in
        node (npu) at (dots.south) [%
        anchor = north,
        yshift = -5pt,
        minimum width = \x2 - \x1,
        minimum height = \y1 - \y2,
        draw,
        signal,
        signal pointer angle = 150,
        signal from = north,
        signal to = nowhere] {NPU};
      \end{tikzpicture}
    \end{column}
  \end{columns}
  \footnotetext[1]<2->{%
    \url{https://cran.r-project.org/doc/manuals/r-release/R-exts.html}}
\end{frame}

\subsection{GPU / NPU support in R has very little baggage}
\begin{frame}{GPU / NPU support in R has very little ``baggage''}
  \begin{columns}[T]
    \begin{column}{.85\framewidth}
      \begin{itemize}[<+->]
      \item Current GPU R packages are {\color{alert}indirect}; %
        simply wrap around python / web interfaces %
        (tensorflow, keras3, ellmer, chattr)
      \item SYCL (compiler), Kokkos and RAJA (libraries) %
        are C++ vendor-neutral heterogeneous computing approaches
      \item But R is a GNU project %
        that will not replace gcc with clang for SYCL support, %
        therefore need a library-approach like Kokkos
      \item<5-> Kokkos needs(?) to be compiled to choose CPUs and GPUs; %
        may need to consider production use of \texttt{spack develop}.
      \item<6-> Use case of Bioconductor R package system
      \end{itemize}
    \end{column}
    \begin{column}{.15\framewidth}
      \begin{tikzpicture}[%
        every node/.style = {
          inner sep = 0pt,
          outer sep = 0pt
        }
        ]
        \visible<4->{
          \node (kokkos) [%
          draw,
          inner sep = 2pt,
          signal,
          signal pointer angle = 150,
          signal from = north,
          signal to = south]
          {\includegraphics[width=\textwidth]{logo-kokkos.png}};
        }
        \path let \p1 = (kokkos.north west), \p2 = (kokkos.south east) in
        node (ffi) at (kokkos.north) [%
        anchor = south,
        yshift = 5pt,
        minimum width = \x2 - \x1,
        minimum height = \y1 - \y2,
        draw,
        signal,
        signal pointer angle = 150,
        signal from = north,
        signal to = south] {FFI};
        \path let \p1 = (kokkos.north west), \p2 = (kokkos.south east) in
        node (r) at (ffi.north) [%
        anchor = south,
        yshift = 5pt,
        minimum width = \x2 - \x1,
        minimum height = \y1 - \y2,
        draw,
        inner sep = 2pt,
        signal,
        signal pointer angle = 150,
        signal from = nowhere,
        signal to = south]
        {\includegraphics[width=.5\textwidth]{logo-r.png}};
        \visible<-3|handout:0>{
          \path let \p1 = (kokkos.north west), \p2 = (kokkos.south east) in
          node (dots) at (kokkos) [%
          minimum width = \x2 - \x1,
          minimum height = \y1 - \y2 - 5pt,
          draw,
          signal,
          signal pointer angle = 150,
          signal from = north,
          signal to = south]
          {\ldots};
        }
        \path let \p1 = (kokkos.north west), \p2 = (kokkos.south east) in
        node (npu) at (kokkos.south) [%
        anchor = north,
        yshift = -5pt,
        minimum width = \x2 - \x1,
        minimum height = \y1 - \y2,
        draw,
        signal,
        signal pointer angle = 150,
        signal from = north,
        signal to = nowhere] {NPU};

        \visible<1|handout:0>{
          \draw [color = alert, line width = 2]
          ($ (ffi.north west) + (-5pt,+5pt) $) --
          ($ (npu.south east) + (+5pt,-5pt) $)
          ($ (ffi.north east) + (+5pt,+5pt) $) --
          ($ (npu.south west) + (-5pt,-5pt) $);
        }
      \end{tikzpicture}
    \end{column}
  \end{columns}
\end{frame}

\subsection{Bioconductor R packages for biological sciences}
\begin{frame}{Bioconductor R packages for biological sciences}
  \begin{columns}[T]
    \begin{column}{.88\framewidth}
      \begin{itemize}[<+->]
      \item Since 2001 by R co-creator Robert Gentleman
        \begin{itemize}[<.->]
        \item $\ge 2000$ packages, %
          $\approx 150k$ unique IP address downloads\footnotemark[1]
        \item Impact of 5000--6000 yearly citations since 2019\footnotemark[1]
        \item First GPU package released this April %
          (RbowtieCuda; not vendor neutral)
        \end{itemize}
      \item CI rebuilds all packages every 48 hours:
        \begin{itemize}[<.->]
        \item Windows Server 2022 (x86\_64) --- \ul{majority of users}
          \begin{itemize}[<3->]
          \item CI Azure builder $\to$ %
            interactive Windows 11 LXC container?
          \end{itemize}
        \item macOS 12, 13 (x86\_64, aarch64)
          \begin{itemize}[<4->]
          \item CI MacStadium.com $\to$ interactive Mac mini? %
            (Apple license limits running macOS on Apple hardware)
          \end{itemize}
        \item Ubuntu Linux 24 (x86\_64, aarch64)
        \end{itemize}
      \end{itemize}
    \end{column}
    \begin{column}{.12\framewidth}
      \begin{tikzpicture}
        \node (bioc)
        {\includegraphics[width=\textwidth]{logo-bioconductor.png}};
        \node (r) at (bioc.south) [anchor = north, yshift = -20pt]
        {\includegraphics[width=\textwidth]{logo-r.png}};
        \draw [-stealth, line width = 2pt] (bioc) -- (r);
      \end{tikzpicture}
    \end{column}
  \end{columns}
  \footnotetext[1]{%
    Bioconductor Annual Report 2024: \url{https://bioconductor.org/about/annual-reports/}}
\end{frame}

\section{Proposal}

\subsection{Proposed R stack for Kokkos NPU support}
\begin{frame}{Proposed R stack for Kokkos NPU support}
  \begin{columns}[T]
    \begin{column}{.85\framewidth}
      \begin{itemize}[<+->]
      \item Bioconductor S4 classes %
        allow strict package interoperability; %
        versioning scheme similar to OS release process
      \item Provide an API to interface with Kokkos C++ using 2~packages
        \begin{enumerate}
        \item<.-> Data \ul{interfaces} and examples
          \begin{itemize}[<.->]
          \item Interfaces for DelayedArray, DelayedTensor S4 classes
          \item Examples using kokkos kernels %
            for \#bio-kokkos relevant algorithms
          \end{itemize}
        \item Expose Kokkos backend \ul{runner} for BiocParallel
          \begin{itemize}
          \item<.-> Existing backends: %
            single CPU, %
            multicore, %
            socket cluster, %
            Redis cluster, %
            MPI cluster, %
            Batchtools (SLURM, etc.)
          \end{itemize}
        \end{enumerate}
      \item \ul{Differs from pykokkos approach} of wrapping kokkos in python;
        BiocKokkos would expose kokkos C++ templates directly.
      \end{itemize}
    \end{column}
    \begin{column}{.15\framewidth}
      \begin{tikzpicture}[%
        every node/.style = {
          inner sep = 0pt,
          outer sep = 0pt
        }
        ]
        \node (kokkos) [%
        draw,
        inner sep = 2pt,
        signal,
        signal pointer angle = 150,
        signal from = north,
        signal to = south]
        {\includegraphics[width=\textwidth]{logo-kokkos.png}};
        \path let \p1 = (kokkos.north west), \p2 = (kokkos.south east) in
        node (ffi) at (kokkos.north) [%
        anchor = south,
        yshift = 5pt,
        minimum width = \x2 - \x1,
        minimum height = \y1 - \y2,
        draw,
        signal,
        signal pointer angle = 150,
        signal from = north,
        signal to = south] {FFI};
        \visible<2->{
          \path let \p1 = (kokkos.north west), \p2 = (kokkos.south east) in
          node (class) at (ffi.north) [%
          anchor = south,
          yshift = 5pt,
          minimum width = \x2 - \x1,
          draw,
          inner sep = 2pt,
          signal,
          signal pointer angle = 150,
          signal from = north,
          signal to = south]
          {BiocKokkos};
        }
        \visible<3->{
          \path let \p1 = (class.north west), \p2 = (class.south east) in
          node (bp) at (class.north) [%
          anchor = south,
          yshift = 5pt,
          minimum width = \x2 - \x1,
          draw,
          inner sep = 2pt,
          signal,
          signal pointer angle = 150,
          signal from = nowhere,
          signal to = south]
          {BiocParallel};
        }
        \visible<-1|handout:0>{
          \path let \p1 = ($ (bp.north west) + (0, 5pt) $),
          \p2 = (class.south east) in
          node (r) at (ffi.north) [%
          anchor = south,
          yshift = 5pt,
          minimum width = \x2 - \x1,
          minimum height = \y1 - \y2,
          draw,
          inner sep = 2pt,
          signal,
          signal pointer angle = 150,
          signal from = nowhere,
          signal to = south]
          {\includegraphics[width=.5\textwidth]{logo-r.png}};
        }
        \path let \p1 = (kokkos.north west), \p2 = (kokkos.south east) in
        node (npu) at (kokkos.south) [%
        anchor = north,
        yshift = -5pt,
        minimum width = \x2 - \x1,
        minimum height = \y1 - \y2,
        draw,
        signal,
        signal pointer angle = 150,
        signal from = north,
        signal to = nowhere] {NPU};
      \end{tikzpicture}
    \end{column}
  \end{columns}
\end{frame}

\section{Wrap up}

\subsection{R already ``bridges'' to a package manager}
\begin{frame}{R already ``bridges'' to a package manager}
  \begin{columns}[T]
    \begin{column}{.75\framewidth}
      \begin{itemize}[<+->]
      \item Bioconductor bioc2u / r2u %
        installs Ubuntu packages (that include dependencies) %
        instead of sources using root privilege escalation and %
        inter-process communication (IPC)
      \item Uses \texttt{bspm R} package (Bridge to System Package Manager)
      \item This is in production\footnotemark[1]
      \item Demonstrates need for a project like spack
        \begin{itemize}
        \item<.-> May be better to lightly couple package source %
          and link to externals with spack for greater flexibility
        \item Support operating systems besides Ubuntu
        \item<5-> \ldots all the other spack advantages
        \end{itemize}
      % \item Traditional R approach would be to lightly wrap %
      %   Kokkos and dependencies as R packages %
      %   exposing libraries and headers
      \end{itemize}
    \end{column}
    \begin{column}{.25\framewidth}
      \begin{tikzpicture}
        \node (bioc)
        {\includegraphics[width=.3\textwidth]{logo-bioconductor.png}};
        \node [anchor = north, yshift = -20pt] (r) at (bioc.south)
        {\includegraphics[width=.3\textwidth]{logo-r.png}};
        \node [anchor = east] (ubuntu) at (bioc.west)
        {\includegraphics[width=.3\textwidth]{logo-ubuntu.png}};
        \visible<4->{
          \node [anchor = west] (spack) at (bioc.east) 
          {\includegraphics[width=.3\textwidth]{logo-spack.pdf}};
          \draw [-stealth, line width = 2pt] (spack) -- (r);
        }
        \draw [-stealth, line width = 2pt] (bioc) -- (r);
        \visible<-3>{
          \draw [-stealth, line width = 2pt] (ubuntu) -- (r);
          \coordinate (mid) at ($ (bioc) !.5! (r) $);
          \draw [color = alert, line width = 2pt]
          (mid) +(-5pt,+5pt) -- +(+5pt,-5pt)
          (mid) +(+5pt,+5pt) -- +(-5pt,-5pt);
          \draw [-stealth, line width = 2pt]
          (bioc) to [out = 120, in = 60, looseness = 2] (ubuntu);
        }
      \end{tikzpicture}
    \end{column}
  \end{columns}
  \footnotetext[1]<3->{\url{https://cloud.r-project.org/bin/linux/ubuntu/}}
\end{frame}

\subsection{3 independent projects}
\begin{frame}{3 independent projects}
  \begin{enumerate}[<+->]
  \item Spack \texttt{bspm} backend to seamlessly install R dependencies %
    with speed up from a build cache
    \begin{itemize}
    \item Then explore ``lightly coupling'' spack with R repositories %
      instead of replacing them % 
      so that R package versions can be upgraded from R instead of spack
    \end{itemize}
  \item spack + kokkos OS support on %
    Windows $\ge 10$\footnotemark[1]
    (Intel Meteor Lake, AMD Ryzen XDNA, Qualcomm Snapdragon X) and %
    macOS $\ge 13.5$\footnotemark[2]
    (Apple silicon)
  \item Bioconductor R packages\footnotemark[3]:
    \begin{description}[Data interface]
    \item[Data interface]<.-> %
      interface with DelayedTensor/Array R packages, Kokkos C++
    \item[Runner] BiocParallel Kokkos backend
    \end{description}
  \end{enumerate}
  \footnotetext[1]<3->{Support for Microsoft Copilot}
  \footnotetext[2]<3->{Support for Apple Metal Machine Learning Explore (MLX)}
  \footnotetext[3]<4->{Will gauge interest at the Bioconductor conference this June}
\end{frame}

\section{Thank you}

{
  \setbeamercolor{normal text}{fg=white}
  \usebeamercolor[fg]{normal text}
  \setbeamercolor{item}{fg=white}
\begin{frame}[noframenumbering, plain]
  \begin{tikzpicture}[remember picture, overlay]
    \node [anchor = north, inner sep = 0pt, outer sep = 0pt]
    at (current page.north) {%
      \includegraphics[width=\paperwidth]{im-hpsf-bg.png}};
    \node at (current page.north) [anchor = north, yshift = -0.7cm] {%
      \LARGE Thank you!};
  \end{tikzpicture}
  \vspace*{0.8cm}
  \begin{columns}[T]
    \begin{column}{.37\framewidth}
      University of Connecticut HPC
      \begin{itemize}
      \item Ed Swindelles
      \item Dave Christianson
      \item Tinyang Xu
      \end{itemize}

      University of Michigan
      \begin{itemize}
      \item Joe Waliga
      \item Denise Kirschner
      \end{itemize}

      University of Pittsburgh
      \begin{itemize}
      \item Tatum Liparulo
      \item Ramakrishna Suresh
      \item Jason Shoemaker
      \end{itemize}
    \end{column}
    \begin{column}{.35\framewidth}
      \vspace{3.5\baselineskip}
      Argonne National Laboratory
      \begin{itemize}
      \item Ray Loy
      \item Sean Koyama
      \item Haritha Som
      \end{itemize}
      Gentoo
      \begin{itemize}
      \item Brandon Mesquita
      \item Roy Bamford
      \item Diego Elio Petten\`o
      \item Donnie Berkholz
      \end{itemize}
    \end{column}
    \begin{column}{.32\framewidth}
      Spack
      \begin{itemize}
      \item Peter Scheibel
      \item Tammy Dalhgren
      \item Adam J. Stewart
      \item Todd Gamblin
      \end{itemize}

      Kokkos, SYCL
      \begin{itemize}
      \item Christian Trott
      \item Thomas Applencourt
      \end{itemize}

      Bioconductor
      \begin{itemize}
      \item Martin Morgan
      \item Luke Tierney
      \item Vince Carey
      \end{itemize}
    \end{column}
  \end{columns}
\end{frame}
}

\end{document}
